\phantomsection
\addcontentsline{lot}{section}{\protect\numberline{}Tabel 36 - Jumlah kematian neonatal dan post neonatal menurut penyebab utama}
\ra{1.3}
% Table generated by Excel2LaTeX from sheet '35'

{\centering
\begin{tabular}{rY{9em}Y{8em}rrrrrrrrrr}
    \multicolumn{13}{l}{Tabel 36}\\
    \multicolumn{13}{c}{JUMLAH KEMATIAN NEONATAL DAN POST NEONATAL MENURUT PENYEBAB UTAMA, KECAMATAN, DAN PUSKESMAS}\\
    \multicolumn{13}{c}{KABUPATEN BELITUNG TIMUR}\\
    \multicolumn{13}{c}{TAHUN \tP}\\
    \toprule
    \multirow{2}[0]{*}{NO} & \multirow{2}[0]{*}{KECAMATAN} & \multirow{2}[0]{*}{PUSKESMAS} & \multicolumn{10}{c}{PENYEBAB KEMATIAN NEONATAL (0-28 HARI)}\\
    \cmidrule(l{2pt}r{2pt}){4-13}
    & & & \multicolumn{1}{c}{\rotatebox{-90}{\makecell[l]{MALFORMASI KONGENITAL, \\DEFORMASI, DAN \\KELAINAN KROMOSOM}}} & \multicolumn{1}{c}{\rotatebox{-90}{\makecell[l]{GANGGUAN TERKAIT \\USIA KEHAMILAN\\ DAN PERTUMBUHAN JANIN}}} & \multicolumn{1}{c}{\rotatebox{-90}{\makecell[l]{TRAUMA \\KELAHIRAN}}} &\multicolumn{1}{c}{\rotatebox{-90}{\makecell[l]{KOMPLIKASI PADA \\SAAT PERSALINAN \\(INTRAPARTUM)}}} &\multicolumn{1}{c}{\rotatebox{-90}{\makecell[l]{KEJANG DAN GANGGUAN\\STATUS SEREBRAL}}} & \multicolumn{1}{c}{\rotatebox{-90}{INFEKSI}} &\multicolumn{1}{c}{\rotatebox{-90}{\makecell[l]{GANGGUAN PERNAPASAN \\DAN KARDIOVASKULAR}}} & \multicolumn{1}{c}{\rotatebox{-90}{\makecell[l]{KONDISI NEONATAL\\ LAINNYA}}} &
    \multicolumn{1}{c}{\rotatebox{-90}{\makecell[l]{BERAT BADAN LAHIR \\RENDAH DAN \\PREMATURITAS}}} & \multicolumn{1}{c}{\rotatebox{-90}{\makecell[l]{KEMATIAN NEONATAL \\DENGAN PENYEBAB YANG \\TIDAK DITENTUKAN}}} \\

    \midrule
    \emph{1} & \emph{2} & \emph{3} & \emph{4} & \emph{5} & \emph{6} & \emph{7} & \emph{8} & \emph{9} & \emph{10} & \emph{11} & \emph{12} & \emph{13}  \\
    \midrule
	1 & Manggar           & Manggar       &  6 & 0 & 0 & 0 & 0 & 0 & 0 & 3 & 0 & 0 \\
	2 & Damar             & Mengkubang    &  2 & 0 & 0 & 1 & 0 & 0 & 0 & 1 & 1 & 0 \\
	3 & Kelapa Kampit     & Kelapa Kampit &  3 & 0 & 0 & 1 & 0 & 0 & 0 & 0 & 0 & 1 \\
	4 & Gantung           & Gantung       &  3 & 0 & 0 & 0 & 1 & 0 & 1 & 0 & 1 & 0 \\
	5 & Simpang Renggiang & Renggiang     &  2 & 0 & 0 & 0 & 0 & 0 & 0 & 0 & 0 & 0 \\
	6 & Simpang Pesak     & Simpang Pesak &  1 & 0 & 0 & 1 & 0 & 0 & 0 & 0 & 0 & 0 \\
	7 & Dendang           & Dendang       &  1 & 0 & 0 & 0 & 0 & 0 & 0 & 0 & 0 & 0 \\
    \midrule
    \multicolumn{3}{l}{JUMLAH KAB.}       & 18 & 0 & 0 & 3 & 1 & 0 & 1 & 4 & 2 & 1 \\
    \bottomrule
\end{tabular}%

}

\vfill
Sumber: Subkoordinator Kesehatan Keluarga dan Gizi\par 
